% ---------------------------------------------------------------------------------------------
% Main text, part 1 (Sections 1-3). Draft written 2026-10-03 from the lean table set (tables_tax.tex)
% and the result CSVs in GACI_FuelShock/. Structure follows Li, Liu, Purevjav and Yang (2019, JEEM).
% Numbers in the text are taken from tables_tax.tex / _res_tax_*.csv; re-check after any re-run of 92-111.
% ---------------------------------------------------------------------------------------------

\section{Introduction}

Air passenger taxes are the most widely used national instrument for pricing the climate cost of flying.
Nine European countries introduced or raised a per-passenger ticket tax between 1998 and 2018, several more
have followed since, and the European Commission's Energy Taxation Directive proposal and the 2027 Dutch
distance-band reform keep the instrument on the agenda. The tax is simple to collect and raises
substantial revenue, but what it does to flying is less clear. A per-passenger levy can reduce emissions
by cutting the number of flights; it can also redistribute flying within the network, because every
European ticket tax exempts transfer passengers and is a larger share of the fare on short, cheap
routes than on long ones. Whether a ticket tax shrinks a country's air network or merely concentrates it
is an empirical question, and the answer determines both the carbon abatement the tax delivers and who
pays for it.

This paper answers that question with the first evaluation of ticket taxes that treats the air network
itself as the outcome. We assemble a monthly panel of scheduled seats and flights for more than 6,000
airports from 1996 to 2024, build flight-stage-resolved CO$_2$ emissions for every departing flight from the
same schedules, and attach to each airport-year a network-centrality summary of its position in the
world air network, the Global Air Connectivity Index (GACI) of \citet{cheung2020}. We then evaluate the
nine European ticket-tax events listed in Table~\ref{tab:events}: the extension of the Danish passenger
tax to domestic flights (1998), the Maltese departure-tax increase (2005), the doubling of UK Air
Passenger Duty (2007), the Dutch ticket tax (2008, repealed in 2009), the Irish Air Travel Tax (2009),
the German and Austrian taxes (2011), and the Norwegian (2016) and Swedish (2018) taxes.

Our empirical framework follows \citet{li2019subway}, who evaluate Beijing's subway expansion with two
complementary designs: a difference-in-differences comparison of locations near and far from each
opening, with a buffer zone to keep misclassified locations out of the control group, and a continuous
network-density measure that captures system-wide effects. We transpose both to the air network. The
difference-in-differences design compares airports in the taxing country with airports in European
countries that have no tax event in the same window, around each tax's announcement. Airports in
untaxed countries that lie within 50 to 300 kilometres of the taxed border enter as a separate exposed
group rather than as controls, because leakage across the border would otherwise contaminate the
comparison, and airports within 50 kilometres are dropped. Event time runs from the announcement, with
the months between announcement and implementation removed, so that anticipation cannot enter either
the pre- or the post-period. The continuous design uses GACI as the network-density analogue: we
estimate how an airport's connectivity moves when its country is taxed, how emissions respond to
connectivity, and how much of the emissions effect runs through the network.

Three results emerge. First, taxes reduce flying at the periphery and leave hubs intact. Across the
eight non-Dutch events, scheduled seats at taxed airports fall by 10 percent in the first tax year when
every airport counts equally, but by 4 percent when airports are weighted by seats. The gap is the
distribution of the effect: airports in the bottom size tercile of the taxing country lose 24 percent of
their seats, 14 percent of their destinations and 4 percentage points of their months with any service,
while large airports lose no seats and increase aircraft size by 2.5 percent. Second, the network
concentrates rather than shrinks. Taxed airports lose 1.6 to 2.3 percent of their connectivity index, 9
percent of their destinations, 19 percent of eigenvector centrality and 43 percent of betweenness
centrality, with the losses again concentrated in small and mid-sized airports; yet country-level
synthetic difference-in-differences estimates on national seat totals and on the seat-weighted national
connectivity index are indistinguishable from zero for every event except Ireland, Malta and Norway.
Third, emissions fall only where flights disappear. Departure CO$_2$ at taxed airports falls by 10
percent unweighted and 2 percent (not significant) seat-weighted; CO$_2$ per seat and per seat-kilometre
do not fall and rise slightly at the hubs that upgauge. An exact decomposition of the CO$_2$ change
attributes all of it to the route-exit margin, with no contribution from flight frequency on surviving
routes or from aircraft and distance. Cross-border leakage is absent in the pooled sample and
measurable only in the Dutch episode, where seats at foreign airports within 150 kilometres rose by 16
percent while the tax was in force.

The paper contributes to three literatures. The first evaluates ticket taxes directly. Synthetic-control,
panel and route-level studies of the German, Austrian, Dutch, Swedish and UK taxes find passenger or
flight reductions of 2 to 12 percent concentrated in low-cost and secondary airports and little or no
effect at hubs \citep{borbely2019,falk2019,bernardo2024flight,oesingmann2022,seetaram2014,gordijn2011}.
We confirm the size gradient on a common footing across nine events, add emissions and an exact
decomposition of how they fall, and show that the gradient is a reallocation of network position, not
just of volume. The second literature evaluates carbon pricing in aviation, chiefly the inclusion of
intra-EEA flights in the EU Emissions Trading System, which lowered emissions on regulated routes by
about 5 percent \citep{fageda2022pricing} and shifted intercontinental connecting passengers toward
hubs outside the system \citep{fageda2025hub}. Ticket taxes differ from allowance prices in two ways
that our results make precise: they are flat per passenger rather than proportional to fuel, so they do
not touch emissions intensity, and they exempt transfers, so they protect the hub function that carbon
pricing is found to erode. The third literature measures air networks with centrality indices
\citep{cheung2020} but has not used them as the outcome of a policy evaluation; the only causal designs
with centrality outcomes concern airline mergers. Our companion work estimates the emissions elasticity
of connectivity at the global level \citep{yoo2026co2} and its trade effects \citep{yoo2026trade}; here
we close the loop by estimating how a specific policy moves connectivity and how the connectivity
change maps into emissions.

The rest of the paper proceeds as in \citet{li2019subway}. Section~\ref{sec:background} describes the
tax events and the data. Section~\ref{sec:strategy} sets out the stacked difference-in-differences,
event-study, synthetic difference-in-differences and decomposition designs. Section~\ref{sec:results}
presents the results for seats, network position and emissions, the mechanism, the national totals,
leakage and the cost per tonne. Section~\ref{sec:conclusion} concludes.

\section{Background and data}
\label{sec:background}

\subsection{Ticket taxes in Europe}

Table~\ref{tab:events} lists the nine events. Every tax is levied per departing passenger on scheduled
commercial flights, is paid by the airline and passed into the fare, and exempts transfer passengers
who connect within a stated time at the taxed airport. Six taxes differentiate by destination distance
in two or three bands (Germany, Austria, Sweden, the Netherlands, Ireland and the United Kingdom); the
Norwegian, Danish and Maltese taxes are flat. Rates at introduction range from about EUR 2 on Irish
short-haul departures to EUR 45 on Dutch and German long-haul departures. Two features of the
institutional setting matter for the design. First, announcement and implementation are separated by
two to eighteen months, and in Germany the tax applied to tickets sold from September 2010 for flights
from January 2011, so demand could react before the effective date. We date events from the
announcement, verified in primary sources (budget speeches, draft bills and official gazettes), and
drop the months between announcement and implementation. Second, the taxes are national, so the
natural control group is airports in European countries without a tax event in the same window, and
the natural threat to that control group is leakage to foreign airports within driving distance of the
taxed country.

\subsection{Airports, schedules and emissions}

The unit of observation is the airport-month. Scheduled seats, flights and seat-kilometres for every
departing flight are taken from the OAG schedule database for January 1996 to June 2024, aggregated to
the airport and month and split into domestic and international departures. CO$_2$ emissions are
computed for each scheduled flight from aircraft-type fuel burn by flight stage (taxi-out, take-off,
climb-out and cruise) and the great-circle distance of the segment, then summed to the airport-month;
the inventory and its validation against EDGAR and national totals are described in
\citet{yoo2026co2}. Attributing cruise emissions to the departure airport avoids double counting and
follows the departure convention of national inventories. We use departure CO$_2$, CO$_2$ per seat and
CO$_2$ per seat-kilometre (emissions intensity), and the average stage length (seat-kilometres per
seat) as outcomes.

\subsection{Network position}

An airport's position in the world air network is measured by the Global Air Connectivity Index
\citep{cheung2020}, recomputed for every year from 1996 to 2024 on the full schedule network of
3,200 to 3,900 active airports. GACI is the first principal component of five standardised centrality
indicators: degree (number of destinations), closeness, eigenvector centrality (connections to
well-connected airports), flow betweenness (transfer flows routed through the airport, computed from an
electrical-circuit representation of the weighted network) and regional importance (intensity of
intra-regional links). The loadings are stable across years and the index is standardised so that
relative scores are comparable over time. We use the index and its degree, eigenvector and betweenness
components as annual outcomes. GACI is a relative measure of position within the year's network, so a
fall at taxed airports means they lost ground relative to other airports, not that the network shrank.

\subsection{Sample and descriptive patterns}

The estimation sample contains airports in 21 European countries for which tax histories were coded
from legal texts (Italy, whose municipal surcharge taxes transfer passengers and has no verified
dates, is excluded from both groups). For each event the treated group is every airport in the taxing
country with scheduled service in the year before the announcement; the control group is every airport
in a coded European country with no tax event within the window; airports in untaxed countries within
50 to 300 kilometres of the nearest treated airport form a border group, and those within 50
kilometres are dropped. Within each taxing country, airports are split into size terciles by seats in
the year before the announcement. The pooled sample has 112,876 airport-months in 33 countries.

A first look at the data follows \citet{li2019subway}. In the twelve months before and after the
effective date, log seats at taxed airports fall from 9.86 to 9.80 while seats at control airports rise
from 9.65 to 9.69. After partialling out airport-by-calendar-month fixed effects, the residual at taxed
airports moves from $+0.034$ to $-0.034$, at control airports from $-0.013$ to $+0.013$, and at border
airports from $-0.015$ to $+0.015$: a raw difference-in-differences of about 7 percent against
controls, with border airports moving with the controls rather than against them. The regressions below
sharpen this comparison.

\section{Empirical strategy}
\label{sec:strategy}

\subsection{Stacked difference-in-differences with a buffer}

For each event $e$ we build a stack of airport-months from 36 months before the announcement $A_e$ to
twelve months after the effective date $E_e$, dropping the months in $[A_e, E_e)$. The main
specification is
\begin{equation}
\ln y_{imt}^{e} = \beta\, T_i^{e} \times \text{Post}_t^{e} + \gamma\, B_i^{e} \times \text{Post}_t^{e}
 + \delta\, X_{c(i)t} + \alpha_{i m}^{e} + \lambda_{t}^{e} + \varepsilon_{imt}^{e},
\label{eq:did}
\end{equation}
where $y$ is seats, flights, seats per flight, CO$_2$, CO$_2$ per seat or CO$_2$ per seat-kilometre at
airport $i$ in calendar month $m$ of month $t$; $T_i^e$ marks airports in the taxing country and
$B_i^e$ airports in the 50--300 km border group; $\text{Post}_t^e$ equals one from $E_e$ onward;
$X_{c(i)t}$ are log GDP and log population of the airport's country; $\alpha_{im}^e$ are airport-by-calendar-month-by-event fixed effects, which absorb each airport's seasonal profile within each stack; and
$\lambda_t^e$ are month-by-event fixed effects. Stacking events with their own fixed effects avoids the
negative weights of two-way fixed-effects estimators under staggered adoption \citep{cengiz2019,goodmanbacon2021}.
Standard errors are clustered by country, the level at which the tax varies. We report equal-weighted
estimates, which describe the typical airport, and seat-weighted estimates, which describe the national
total, and we interact the treatment with the within-country size tercile.

The buffer-and-ring treatment of foreign airports is the analogue of the 2--20 km buffer in
\citet{li2019subway}. Airports just across the border are the ones most likely to gain passengers who
drive to avoid the tax; treating them as controls would bias the estimate toward larger losses, and
dropping them would discard the leakage margin. Entering them as their own group identifies the leakage
coefficient $\gamma$ directly and keeps the control group clean \citep{butts2023}. The Dutch episode,
in which the tax was introduced in July 2008 and set to zero in July 2009, is analysed separately as an
introduction followed by a removal and excluded from the pooled stacks.

\subsection{Event time and sensitivity to pre-trends}

The identifying assumption is that, absent the tax, seats at taxed airports would have followed the
control path. We test it with event-time versions of \eqref{eq:did} in twelve-month bins so that
seasonality cannot enter: two pre-announcement bins ($[A{-}36, A{-}25]$ and $[A{-}24, A{-}13]$),
the reference bin $[A{-}12, A{-}1]$, and post bins $[E, E{+}11]$ and $[E{+}12, E{+}23]$, pooled and event by
event. For annual outcomes (GACI and its components, and the decomposition of CO$_2$) the stacks run
from three years before to one year after the effective year with airport-by-event and year-by-event
fixed effects. We complement the pre-trend tests with the sensitivity analysis of \citet{rambachan2023},
which reports the set of post-period estimates consistent with post-treatment violations of parallel
trends up to a multiple $\bar M$ of the largest pre-period violation.

\subsection{Synthetic difference-in-differences for national totals}

The airport-level estimates describe the distribution of the effect across airports. To ask whether a
country's total flying and its national network position changed, we estimate the synthetic
difference-in-differences of \citet{arkhangelsky2021} event by event on the country-by-month log of
national seats and on the country-by-year seat-weighted mean of log GACI. Donors are the coded European
countries without a tax event in the window; unit weights match the treated country's pre-announcement
path and time weights match the pre-period to the post-period; inference uses placebo estimates that
treat each donor in turn.

\subsection{Decomposing the emissions effect and the network channel}

Annual departure CO$_2$ at an airport satisfies the identity
\begin{equation}
\ln \text{CO}_2 = \ln(\text{destinations}) + \ln(\text{flights per destination}) + \ln(\text{CO}_2 \text{ per flight}),
\label{eq:decomp}
\end{equation}
and CO$_2$ per flight further equals seats per flight times stage length times emissions per
seat-kilometre. Estimating \eqref{eq:did} on each term splits the emissions effect into a network
margin (routes dropped), a frequency margin (flights on surviving routes) and an aircraft-and-distance
margin, with coefficients that add up by construction. To relate the tax effect to connectivity, we
estimate the within-airport elasticity of CO$_2$ to GACI with airport and country-by-year fixed
effects, and apply the decomposition of \citet{gelbach2016} to split the tax effect into the part that
runs through the change in GACI and the part that does not. We treat this as a descriptive channel
analysis rather than a causal mediation estimate: an instrument for connectivity that is excludable from
the emissions equation and strong in the tax stacks is not available (Appendix), so we do not claim
that the network share identifies a structural indirect effect.

\subsection{Why not a country panel}

A natural alternative is a country-year panel of national emissions on a tax indicator with country and
year fixed effects. We report in the Appendix that such a regression is not robust: taxing countries are
concentrated in Northern and Central Europe, whose air markets grew more slowly after 2010 than the
Southern and Eastern European leisure markets that make up the never-taxed group, for reasons unrelated
to the tax, and the sign of the country-level coefficient depends on whether country trends are
included. The stacked design with announcement timing, a short window, airport-by-calendar-month fixed
effects and the buffer group is what removes this confounding, which is why we follow the local
comparison of \citet{li2019subway} rather than a cross-country regression.
